%%%%%%%%%%%%%%%%%%%%%%%%%
\spellentry{Passwall}
%%%%%%%%%%%%%%%%%%%%%%%%%

\tagline{Transmutation}
\begin{description*}
\item[Level:] Sor/Wiz 5
\item[Casting Time:] 1 standard action
\item[Range:] Touch
\item[Effect:] 5 ft. by 8 ft. opening, 10 ft. deep plus 5 ft. deep per three
\item[Duration:] 1 hour/level (D)
\item[Saving Throw:] None
\item[Spell Resistance:] No
\item[Components:] V, S, M
\item[Arcane Material Component:] A pinch of sesame seeds.
\end{description*}

additional levels
You create a passage through wooden, plaster, or stone walls, but not through metal
or other harder materials. The passage is 10 feet deep plus an additional 5 feet
deep per three caster levels above 9th (15 feet at 12th, 20 feet at 15th, and a
maximum of 25 feet deep at 18th level). If the wall's thickness is more than the
depth of the passage created, then a single \textit{passwall} simply makes a niche
or short tunnel. Several \textit{passwall} spells can then form a continuing passage
to breach very thick walls. When \textit{passwall} ends, creatures within the passage
are ejected out the nearest exit. If someone dispels the \textit{passwall} or you
dismiss it, creatures in the passage are ejected out the far exit, if there is
one, or out the sole exit if there is only one.